Tutores: Francisco Herrera e Iván Sevillano

0. Índice

\hyperref[entrenamiento-de-distintos-modelos]{\textbf{1. Entrenamiento
de distintos modelos 3}}

\begin{quote}
\hyperref[base-de-dato-de-entrenamientoid]{1.1. Base de dato de
entrenamiento(ID) 3}

\hyperref[dataloader]{1.2. DataLoader 4}

\hyperref[learner]{1.3. Learner 4}

\hyperref[anuxe1lisis-de-resultados]{1.4. Análisis de resultados 5}

\hyperref[gruxe1ficos-y-resultados-del-entrenamiento-de-cada-modelo]{1.5.
Gráficos y resultados del entrenamiento de cada modelo 6}

\hyperref[modelo-primero]{1.5.1. Modelo Primero: 6}

\hyperref[modelo-2]{1.5.2. Modelo 2 7}

\hyperref[modelo-3]{1.5.3. Modelo 3: 7}

\hyperref[modelo-4]{1.5.4. Modelo 4: 8}

\hyperref[comparativa-de-los-modelos-obtenidos-y-sus-resultados-asociados.]{1.6.
Comparativa de los modelos obtenidos y sus resultados asociados. 9}
\end{quote}

\hyperref[out-of-distribution]{\textbf{2. Out of Distribution 9}}

\begin{quote}
\hyperref[conjunto-de-datos]{\textbf{2.1. Conjunto de Datos 10}}

\hyperref[muxe9tricas]{\textbf{2.2. Métricas 13}}

\hyperref[muxe9todos-de-detecciuxf3n-ood-post-hoc]{2.3. Métodos de
detección OOD post hoc 13}

\hyperref[grupo-1]{2.3.1. Grupo 1 13}

\hyperref[grupo-2-comparan-con-la-informaciuxf3n-de-las-features-de-entrenamiento]{2.3.2.
Grupo 2: comparan con la información de las features de entrenamiento
14}

\hyperref[grupo-3-muxe9todos-basados-en-la-modificaciuxf3n-de-activaciones]{2.3.3.
Grupo 3: métodos basados en la modificación de activaciones 14}

\hyperref[resultados-obtenidos]{2.3.3. Resultados obtenidos 15}

\hyperref[anuxe1lisis-de-los-resultados-de-detecciuxf3n-ood]{2.3.4.
Análisis de los resultados de detección OOD 15}

\hyperref[stood-x]{2.3.5. STOOD-X 16}
\end{quote}

\section{Entrenamiento de distintos
modelos}\label{entrenamiento-de-distintos-modelos}

\subsection{1.1. Base de dato de
entrenamiento(ID)}\label{base-de-dato-de-entrenamientoid}

Con punto de partida contábamos con una base de datos de imágenes
médicas clasificadas según el tipo de enfermad que presentaban. La base
contaba con las siguientes clases: Adenocarcinoma\_ductal,
Glandulas\_no\_neoplasicas, Gleason\_4\_fusionadas
,Gleason\_5\_comedonecrosis, No\_Cancer ,Carcinoma\_intraductal
Gleason\_3 , Gleason\_4\_glomeruloide, Gleason\_5\_no\_glandulas,
PIN\_de\_alto\_grado, Estroma , Gleason\_4\_cribiformes,
Gleason\_4\_pobremente\_formadas, Inflamatorio. Inicialmente el
repositorio cuenta de partida con dos subcarpetas que son train y val,
con un total de 126192 imágenes, cuya distribución de imágenes por cada
clase es la siguiente:

\includegraphics[width=4.36979in,height=3.27371in]{notas/media/image3.png}

Contado con un total de 90800 imágenes para train, 22691 para validation
y 12618 para test.

El primer reto al que nos enfrentamos es el de poder crear un modelo de
aprendizaje supervisado que dado una imágen de nuestro conjunto de datos
sea capaz de clasificarla en una de las respectivas clases con la mayor
precisión posible.

Para comenzar se necesita un conjunto train, val y test. Así pues dada
estructura inicial del conjunto de datos, se optó por utilizar el
directorio val, que presentaba una menor cantidad de imágenes en
comparación con train para realizar el testing, y a su vez de train se
separaría en dos carpetas, una mayor para el entrenamiento y una menor
para la validación.

\subsection{1.2. DataLoader}\label{dataloader}

Para una primera aproximación del modelo hemos decidido crear un
conjunto de imágenes menor apartir del original, el cual sirve para
pruebas de testeo y comprobar el funcionamiento correcto de nuestro
modelo antes de pasarlo a entrar con el conjunto de datos original. Este
al igual que el conjunto de datos original cuenta con tres carpetas,
inicialmente se seleccionaron 10 imágenes de cada carpeta de la
estructura original del dataset para crear uno nuevo pero con un número
mucho más reducido de imágenes. Posteriormente se aplico el mismo
procedimiento de separar train en dos nuevas carpetas de train y val, y
el val original pasó a ser test. Debido a esto el resultado final ha
sido una estructura de tres carpetas train con 139 imágenes, val con 55
y test con 167.

\includegraphics[width=3.5464in,height=2.6576in]{notas/media/image18.png}

\subsection{1.3. Learner}\label{learner}

En cuanto al optimizador vamos a usar el AdamW {[}1{]} pues se trata de
una versión mejorada del optimizador Adam con la mejora añadida de que
cuenta con Weight Decay. En la implementación de FastAi es cierto que el
nombre asociado es el de Adam, pero es debido a una decisión de FastAi
de adoptar la versión mejorarda del optimizado bajo el nombre de la
original.

Debido a la complejidad del problema no tardó en saltar el problema de
memoria CUDA out of Memory Error. Para plantearnos manejar este problema
sin alterar la arquitectura del modelo {[}3{]} se optó por reducir el
tamaño del batch size a la mitad, pasando así entonces de 128 imágenes
por lote a 64. Las otras medidas tomadas fueron adaptar el learner para
trabajar con precisión de 16 bits, ahorrando así memoria, liberar la
caché de la memoria al iniciar la ejecución, y por último se ha reducido
el número de workers de 8 a 4, planteando la posiblidad de que la
creación de hebras separadas para el tratamiento del problema pudiera
estar infiriendo en una mayor lentitud de ejecución del modelo.

Llegados a este punto solo queda probar con diversos modelos ya
preentrenados, hacerles un fine tuning y analizar cual genera mejor
resultados y el que será por tanto con el que sigamos trabajando a lo
largo del proyecto:

\subsection{1.4. Análisis de
resultados}\label{anuxe1lisis-de-resultados}

{\def\LTcaptype{none} % do not increment counter
\begin{longtable}[]{@{}ll@{}}
\toprule\noalign{}
\textbf{RESUMEN} & test\_accuracy \\
vit\_base\_patch16\_224 & 0,9525281190872192 \\
resnet50 & 0.8963385820388794 \\
convnext\_small & 0.9667934775352478 \\
tf\_efficientnetv2\_s & 0.9537961483001709 \\
swin\_small\_patch4\_window7\_224 & 0.9686955213546753 \\
\midrule\noalign{}
\endhead
\bottomrule\noalign{}
\endlastfoot
\end{longtable}
}

El modelo Swin Transformer obtiene la mayor accuracy (96,87\%), seguido
muy de cerca por ConvNeXt (96,68\%). Dada la escasa diferencia entre
ambos, cualquiera de los dos sería una base adecuada; se selecciona Swin
por ser el que mejor resultado ha obtenido.''

La superioridad de esta arquitectura radica en tres pilares
fundamentales:

\begin{enumerate}
\def\labelenumi{\arabic{enumi}.}
\item
  \textbf{Jerarquización y Ventanas Desplazadas (Shifted Windows):} A
  diferencia de los Vision Transformers (ViT) estándar, que operan con
  parches globales, el Swin Transformer utiliza un mecanismo de ventanas
  jerárquicas. Esto permite capturar relaciones espaciales tanto a nivel
  local (textura, morfología de tejidos) como global (anatomía
  estructural de la lesión).
\item
  \textbf{Eficiencia en el Contexto Local:} La capacidad de
  contextualizar parches mediante window7 permite que el modelo se
  enfoque en características relevantes sin verse afectado por el ruido
  de fondo, algo que los modelos de atención global (como
  vit\_base\_patch16) pueden ignorar al tratar todas las regiones por
  igual.
\item
  \textbf{Inductividad de la Imagen:} Al igual que ConvNeXt, el Swin
  Transformer integra principios de las redes convolucionales (sesgo
  inductivo), lo que facilita el aprendizaje de características
  jerárquicas desde resoluciones finas a gruesas. Esto permite una
  convergencia más eficiente y un mejor aprovechamiento de la
  información contenida en las imágenes médicas.
\end{enumerate}

La diferencia obtenida mínima entre el \emph{Swin Transformer} y
\emph{ConvNeXt} (0.0019) indica que ambas arquitecturas están alcanzando
un techo de rendimiento cercano al límite de información útil extraíble
del conjunto de datos.

\begin{itemize}
\item
  \textbf{Complementariedad de enfoques:} El hecho de que modelos
  basados en atención jerárquica y arquitecturas convolucionales
  modernas compitan por el primer puesto sugiere que las características
  más relevantes para el diagnóstico están siendo bien capturadas por
  ambos paradigmas.
\end{itemize}

Por tanto el modelo con el que nos vamos a quedar para proseguir en
nuestro estudio es el swin\_small\_patch4\_window7\_224.

\subsubsection{}\label{section}

\subsection{1.5. Gráficos y resultados del entrenamiento de cada
modelo}\label{gruxe1ficos-y-resultados-del-entrenamiento-de-cada-modelo}

\subsubsection{1.5.1. Modelo Primero:}\label{modelo-primero}

\includegraphics[width=7.45869in,height=1.89565in]{notas/media/image4.png}

\includegraphics[width=3.52826in,height=3.51984in]{notas/media/image5.png}

Test Loss: 0.6787635087966919, Test Accuracy: 0.9525281190872192

\subsubsection[1.5.2. Modelo 2]{\texorpdfstring{1.5.2. Modelo
2\protect\includegraphics[width=7.44094in,height=1.86024in]{notas/media/image15.png}}{1.5.2. Modelo 2}}\label{modelo-2}

\includegraphics[width=3.52756in,height=3.5236in]{notas/media/image1.png}

Test Loss: 0.8891878724098206, Test Accuracy: 0.8963385820388794

\subsubsection[1.5.3. Modelo 3:]{\texorpdfstring{1.5.3. Modelo
3:\protect\includegraphics[width=7.44094in,height=1.93465in]{notas/media/image7.png}}{1.5.3. Modelo 3:}}\label{modelo-3}

\includegraphics[width=3.52756in,height=3.5236in]{notas/media/image16.png}

Test Loss: 0.6424340009689331, Test Accuracy: 0.9667934775352478

\subsubsection[1.5.4. Modelo 4:]{\texorpdfstring{1.5.4. Modelo
4:\protect\includegraphics[width=7.44094in,height=1.93465in]{notas/media/image2.png}}{1.5.4. Modelo 4:}}\label{modelo-4}

\includegraphics[width=3.52756in,height=3.5236in]{notas/media/image14.png}

\subsection{\texorpdfstring{1.6. Comparativa de los modelos obtenidos y
sus resultados asociados.
}{1.6. Comparativa de los modelos obtenidos y sus resultados asociados. }}\label{comparativa-de-los-modelos-obtenidos-y-sus-resultados-asociados.}

\section{2. Out of Distribution}\label{out-of-distribution}

Vamos a usar el repositorio alojado en este
\href{https://github.com/Jingkang50/OpenOOD}{\ul{github}} para continuar
con nuestro análisis de resultados, particularmente en este
\href{https://colab.research.google.com/drive/1tvTpCM1_ju82Yygu40fy7Lc0L1YrlkQF?usp=sharing\#scrollTo=3jf2U83Apcru}{\ul{enlace}}
se encuentra una explicación de uso del material:

\includegraphics[width=4.81352in,height=2.88651in]{notas/media/image9.png}

\subsection{2.1. Conjunto de Datos}\label{conjunto-de-datos}

{\def\LTcaptype{none} % do not increment counter
\begin{longtable}[]{@{}ll@{}}
\toprule\noalign{}
\textbf{Conjunto} & \textbf{Imágenes (14 clases)} \\
Train ID & 90.800 \\
Validación ID & 22.691 \\
Test ID & 12.618 \\
\midrule\noalign{}
\endhead
\bottomrule\noalign{}
\endlastfoot
\end{longtable}
}

{\def\LTcaptype{none} % do not increment counter
\begin{longtable}[]{@{}
  >{\raggedright\arraybackslash}p{(\linewidth - 6\tabcolsep) * \real{0.2500}}
  >{\raggedright\arraybackslash}p{(\linewidth - 6\tabcolsep) * \real{0.2500}}
  >{\raggedright\arraybackslash}p{(\linewidth - 6\tabcolsep) * \real{0.2500}}
  >{\raggedright\arraybackslash}p{(\linewidth - 6\tabcolsep) * \real{0.2500}}@{}}
\toprule\noalign{}
\textbf{Conjunto} & \textbf{Papel} & \textbf{¿Qué cambia respecto al
ID?} & \textbf{Objetivo} \\
NCT-CRC: TUM + NORM (23080)

(A) & Near OOD & Otro órgano (colon, tejido glandular) y otro centro;
misma escala & Buscamos ver si un tejido muy parecido al original es
capaz de rechazarlo \\
NCT-CRC: STR, LYM, MUS (por separado)

(B) & Control interno & Otro centro; tejidos parecidos a clases de
entrenamiento (Estroma e inflamatorio) & Tejidos nuevamente parecidos a
ID, determinamos si descarta por el centro por el tejido \\
PANDA(10000)

(C) & Covariate shift & Mismo órgano (próstata), otro centro y escáner;
menor resolución (128 px) & En este caso al repetir el organo lo único
que cambia que afecta es el covariate shift por lo que una detección OOD
reflejaría su discrminación por este motivo \\
Control de resolución (12618) & Control & Las mismas imágenes del test
ID, reducidas a 128 px y ampliadas a 224 & Buscamos ver si se están
clasificando imágenes en OOD solo por la resolución \\
DTD / Textures & Far -OOD & No es tejido: texturas genéricas &
Clasificación sencilla, ha de clasificarla directamente como OOD \\
NCT-CRC:ADI, BACK, DEB, MUC & ANEXO & Grasa, fondo, restos, moco & \\
\midrule\noalign{}
\endhead
\bottomrule\noalign{}
\endlastfoot
\end{longtable}
}

\textbf{ID (información extraída del paper de Iván Sevillano sobre
Cáncer de Próstata)}

El muestreo tisular fueron realizados por residentes de anatomía
patológica o por patólogos titulados. Las piezas se procesaron en un
procesador automático de tejidos, se deshidrataron y se incluyeron en
parafina. Se obtuvieron cortes de 3--5 μm de grosor con un micrótomo,
que se montaron en portaobjetos de vidrio, se desparafinaron y se
tiñeron con hematoxilina y eosina (H\&E).

Las preparaciones teñidas con H\&E se digitalizaron con un escáner de
preparaciones completas Aperio/Leica AT2 a 20× aumentos, con una
resolución de escaneo de 50 000 píxeles por pulgada (0,5 μm/píxel).
Todas las imágenes fueron segmentadas manualmente por patólogos
experimentados con un amplio conocimiento de la histología del cáncer de
próstata. Las regiones anotadas se clasificaron en 14 categorías
histológicas distintas que reflejan patrones morfológicos reconocidos en
la práctica clínica.

Cada WSI anotada se dividió en parches cuadrados de tamaño H × H = 224 ×
224 píxeles, extraídos con un paso (\emph{stride}) de p = 256 píxeles.
Un parche se asignaba a una clase concreta si más del 75 \% de sus
píxeles pertenecían a una región anotada como dicha clase; en caso
contrario, se descartaba.

\textbf{NCT-CRC-HE-100K-NONORM}

\begin{itemize}
\item
  Este es un conjunto de 100.000 parches de imágenes no superpuestas de
  imágenes histológicas teñidas con hematoxilina y eosina (H\&E) de
  cáncer colorrectal humano (CRC) y tejido normal.
\item
  Todas las imágenes son 224x224 píxeles (px) a 0,5 micras por píxel
  (MPP).Se utiliza la versión NONORM, generada a partir de los mismos
  datos originales pero sin normalización de color, para no introducir
  un cambio artificial en la
  distribución.''(http://ieeexplore.ieee.org/abstract/document/5193250/,
  \href{https://doi.org/10.1109/ISBI.2009.5193250}{\ul{DOI
  10.1109/ISBI.2009.}}5193250).
\item
  Las clases de tejido son: tejido adiposo (ADI), fondo (BACK), desechos
  (DEB), linfocitos (LYM), moco (MUC), músculo liso (MUS), mucosa normal
  del colon (NORM), estroma asociado al cáncer (STR), epitelio de
  adenocarcinoma colorrectal (TUM).
\item
  Estas imágenes se extrajeron manualmente de N=86 H\&E portaobjetos de
  tejido de cáncer humano teñidos de muestras de parafina embebidas en
  formalina (FFPE) del NCT Biobank (Centro Nacional de Enfermedades
  Tumorales, Heidelberg, Alemania) y el archivo de patología UMM
  (University Medical Center Mannheim, Mannheim, Alemania). Las muestras
  de tejido contenían portaobjetos tumorales primarios de CRC y tejido
  tumoral de metástasis hepáticas de CRC; las clases de tejido normal se
  aumentaron con regiones no tumorales de la muestra de gastrectomía
  para aumentar la variabilidad.
\end{itemize}

\begin{enumerate}
\def\labelenumi{(\Alph{enumi})}
\item
  TUM y NORM hacen referencia epitelio de adenocarcinoma colorrectal y a
  mucosa normal del colon
\end{enumerate}

\begin{quote}
\includegraphics[width=1.83482in,height=1.84167in]{notas/media/image11.png}\includegraphics[width=1.83232in,height=1.8711in]{notas/media/image10.png}

Comparativa de la adenocarcinoma colorrectal versus adenocarcinoma
ductal del ID
\end{quote}

\begin{enumerate}
\def\labelenumi{(\Alph{enumi})}
\setcounter{enumi}{1}
\item
  STR, LYM y MUS hacen referencia a estroma asociado al cáncer,
  linfocitos y músculo liso
\end{enumerate}

\begin{quote}
\includegraphics[width=2.15208in,height=2.13626in]{notas/media/image12.png}\includegraphics[width=2.1625in,height=2.11444in]{notas/media/image13.png}

Ejemplos de imágenes de STR y LYM

(C) Panda

\includegraphics[width=1.59375in,height=1.63542in]{notas/media/image6.png}\includegraphics[width=1.61875in,height=1.63065in]{notas/media/image17.png}

Panda versus nuestras imagenes de ID
\end{quote}

\href{https://zenodo.org/records/1214456}{\ul{Enlace al conjunto de
datos}}

\textbf{Covariate Shift--- PANDA (próstata de otro centro):} Biopsias
del mismo órgano (próstata) y mismo tipo de tinción, pero de escáneres y
hospitales distintos (Karolinska Institute / Radboud University).

\href{https://www.kaggle.com/c/prostate-cancer-grade-assessment}{\ul{https://www.kaggle.com/c/prostate-cancer-grade-assessment}}

\href{https://www.kaggle.com/iafoss/panda-16x128x128-tiles-data}{\ul{https://www.kaggle.com/iafoss/panda-16x128x128-tiles-data}}

\subsection{2.2. Métricas}\label{muxe9tricas}

Hay tres métricas que miden lo bien que mi clasificador separa ID de
OOD. El inconveniente está en que no mide si la clasificación de las 14
clases es correcta. Para ello conviene mejor usar CCR@FPRx. Al final
estas métricas se comportan como modelos de clasificación binaria en ID
y OOD:

\begin{itemize}
\item
  \textbf{AUROC}: representa la probabilidad de que el modelo asigne una
  puntuación o probabilidad más alta a un ejemplo positivo elegido al
  azar que a uno negativo. \textbf{Rango de valores:} Va de 0 a 1.0
  representado 0.5 que el modelo no tiene capacidad discriminatoria
  (equivale a adivinar al azar), \textbf{1.0} que el modelo tiene una
  clasificación perfecta y aquellos casos con \emph{valores inferiores a
  0.5 i}ndican que el modelo invierte las clases (si se invierte la
  predicción, el valor pasa a ser 1 - valor). En este caso también la
  clase positiva es la OOD.
\item
  \textbf{FPR@95:} consiste en ajustar el umbral que delimita un caso
  como ID de forma que dicho umbral detecte el 95\% de las OOD.
\item
  AUPR: esta métrica va acompañada de AUROC, ya que AUROC no es sensible
  a desbalance entre clases mientras que AUPR sí. Principalmente se basa
  en dos conceptos Precisión y Recall. De esta métrica obtenemos dos
  versiones AUPR-In (mide lo bien que se recuperan imágenes de próstata
  sin que se cuelen OOD) y AUPR-Out (mide lo bien que se encuentran las
  OOD sin falsas alarmas sobre las próstata), en el primer caso los
  casos positivos en ID y por eso mismo nos interesa un valor alto y en
  el segundo el caso positivo es el OOD por lo que interesa un valor más
  bajo.
\item
  \textbf{CCR@FPRx}: Se fija el umbral de modo que solo se acepte un x
  \% de las imágenes OOD, y se mide la proporción de imágenes del test
  ID que son aceptadas y además correctamente clasificadas. Se calcula
  solo sobre el test ID; su máximo es la exactitud del clasificador.
\end{itemize}

\subsection{2.3. Métodos de detección OOD post
hoc}\label{muxe9todos-de-detecciuxf3n-ood-post-hoc}

\subsubsection{2.3.1. Grupo 1}\label{grupo-1}

Grupo de los logits (números producidos por la red justo al final, antes
de convertirlos en probabilidades):

\begin{itemize}
\item
  \textbf{MSP (Maximum Softmax Probability)}: la probabilidad de la
  clase ganadora tras el softmax.
\item
  \textbf{MLS (MaxLogit)}: lo mismo que MSP pero con el logit máximo sin
  pasar por softmax. La ventaja que ofrece este método es que el softmax
  normaliza y pierde información sobre la magnitud absoluta; el logit
  bruto la conserva.
\item
  \textbf{EBO (Energy-Based OOD)}: el logaritmo de la suma de las
  exponenciales de los argumentos de todos los logits. Funciona como una
  versión "suave" y diferenciable de la función max(x).
\end{itemize}

\subsubsection{2.3.2. Grupo 2: comparan con la información de las
features de
entrenamiento}\label{grupo-2-comparan-con-la-informaciuxf3n-de-las-features-de-entrenamiento}

\begin{itemize}
\item
  \textbf{KNN}: guarda las features de entrenamiento (normalizadas) y,
  para cada imagen de test, mide la distancia a su k-ésimo vecino más
  cercano. Si está lejos de todo lo visto, es OOD.
\item
  \textbf{MDS (Mahalanobis)}: calcula la media de las features de cada
  clase y una covarianza compartida. La puntuación es la distancia de
  Mahalanobis a la media de la clase más cercana
\item
  \textbf{RMDS (Relative Mahalanobis)}: igual que MDS, pero resta la
  distancia a una gaussiana global de todo el dataset. Así corrige el
  hecho de que algunas direcciones del espacio de features son ``fondo''
  común a todas las imágenes y no información de clase. Suele funcionar
  mejor en near-OOD.
\item
  \textbf{ViM (Virtual-logit Matching)}: busca las direcciones
  principales del espacio de features de entrenamiento (análisis de
  componentes principales). Lo que queda fuera de ese subespacio
  principal (el residuo) se considera señal de OOD. Combina la norma de
  ese residuo con la energía de los logits.
\item
  \textbf{SHE (Simplified Hopfield Energy)}: en el setup() pasa el train
  por la red y, para cada clase, calcula la media de las features sólo
  de las imágenes que el modelo clasifica correctamente (etiqueta real =
  predicción). Eso le da 14 ``patrones'' de referencia. En test, mira
  qué clase predice el modelo para la imagen y compara su feature con el
  patrón de esa clase (por producto escalar, distancia euclídea o
  coseno, según la configuración).
\end{itemize}

\subsubsection{2.3.3. Grupo 3: métodos basados en la modificación de
activaciones}\label{grupo-3-muxe9todos-basados-en-la-modificaciuxf3n-de-activaciones}

Estos métodos modifican las activaciones o los pesos de la última capa
del modelo antes de calcular la puntuación, con el objetivo de atenuar
las señales anómalas que producen las imágenes OOD. Los tres calculan
finalmente una puntuación de tipo Energy sobre los logits resultantes y
fueron diseñados y evaluados originalmente sobre redes convolucionales
con activaciones ReLU.

\begin{itemize}
\item
  \textbf{ReAct (Rectified Activations)} (Sun et al., 2021): parte de la
  observación de que las imágenes OOD tienden a producir activaciones
  anormalmente altas en algunas unidades de la penúltima capa, lo que
  sus autores atribuyen al uso de estadísticas de BatchNorm estimadas
  con datos ID y aplicadas sin cambios a datos OOD. Para atenuarlo,
  recorta cada activación del vector de características por encima de un
  umbral c, h̃ = mín(h, c), y calcula la puntuación Energy sobre los
  logits obtenidos a partir de las características recortadas. El umbral
  c se fija como el percentil 90 de las activaciones en el conjunto de
  validación ID.
\item
  \textbf{ASH (Activation Shaping)} (Djurisic et al., 2023): actúa sobre
  el vector de características de cada imagen de forma individual.
  Conserva solo el 10 \% de las activaciones de mayor valor (percentil
  90), pone a cero el resto y ajusta las conservadas para compensar la
  información eliminada. En la variante empleada (ASH-B), todas las
  activaciones conservadas toman el mismo valor, igual a la suma total
  de las activaciones de la imagen dividida entre el número de
  activaciones conservadas. Después se calcula la puntuación Energy. Se
  aplica sobre las 768 características agregadas tras el \emph{pooling},
  como en la implementación de OpenOOD.
\item
  \textbf{DICE (Directed Sparsification)} (Sun y Li, 2022): actúa sobre
  los pesos de la última capa lineal en lugar de sobre las activaciones.
  Para cada conexión entre una característica y una clase, calcula su
  contribución media en el conjunto de entrenamiento ID (el peso
  multiplicado por la activación media de esa característica), conserva
  solo las conexiones de mayor contribución y pone a cero el resto
  (esparsidad del 90 \%). Al eliminar las conexiones poco relevantes
  para los datos ID, reduce la varianza de la salida en las imágenes OOD
  y mejora su separación. La puntuación es la Energy de los logits
  calculados con la capa podada.
\end{itemize}

\subsubsection{2.3.3. Resultados obtenidos}\label{resultados-obtenidos}

pdf - tabla\_principal\_ood.tex

\subsubsection{2.3.4. Análisis de los resultados de detección
OOD}\label{anuxe1lisis-de-los-resultados-de-detecciuxf3n-ood}

Los resultados muestran una separación clara entre familias de métodos.
Los métodos que operan sobre el espacio de características (ViM, MDS,
KNN, RMDS y STOOD-X) superan de forma sistemática a los basados
únicamente en los logits, en todos los conjuntos evaluados. En el
near-OOD principal (NCT-CRC TUM + NORM), los primeros alcanzan un AUROC
de entre 0,953 y 0,969, frente a 0,821--0,876 de MSP, MaxLogit y Energy.
La diferencia es aún mayor en el punto de trabajo: el FPR@95 baja de
0,879 (MSP) a valores entre 0,128 y 0,173, y el CCR@FPR5 sube de 0,118 a
más de 0,82. Esto indica que el vector de características de 768
dimensiones contiene información sobre la pertenencia a la distribución
que se pierde al proyectarlo sobre los 14 logits.

\textbf{Near-OOD principal.} ViM (0,969 {[}0,968; 0,971{]}) y MDS
(0,968) obtienen el mejor AUROC, mientras que RMDS logra el mejor FPR@95
(0,128) y el mejor CCR@FPR5 (0,868), con un AUROC significativamente
inferior (0,954). Ambos resultados no son contradictorios: el AUROC
resume la separación en todos los umbrales, mientras que FPR@95 y
CCR@FPR5 evalúan la zona de umbral exigente, que es la relevante en un
contexto clínico. Con el umbral que deja pasar solo el 5 \% de las
imágenes de colon, RMDS acepta y clasifica correctamente el 86,8 \% del
test ID, frente a un máximo posible de 96,9 \%. No obstante, RMDS
presenta un fallo marcado en la clase LYM (FPR@95 = 0,992; CCR@FPR5 =
0,008): una parte de los linfocitos de colon recibe puntuaciones
superiores a las de las propias imágenes ID, por lo que ningún umbral
permite rechazarlos sin descartar casi todo el ID.

\textbf{Control interno: centro frente a tejido.} Las clases afines de
NCT-CRC permiten separar el efecto del cambio de centro del efecto del
tipo de tejido. En ViM, MDS y KNN, los linfocitos (LYM), el tejido más
similar entre órganos, se detectan con un AUROC (0,948--0,969)
prácticamente igual al del tejido glandular de colon (TUM: 0,952--0,966;
NORM: 0,963--0,975). Por tanto, estos detectores responden
principalmente al \textbf{cambio de centro} (laboratorio, tinción,
escáner) y no a la naturaleza del tejido. Este resultado es coherente
con lo señalado por Yang et al. (2024): los detectores OOD tienden a ser
más sensibles al desplazamiento de covariables que al cambio semántico.

\textbf{Cambio de centro y resolución.} PANDA, compuesto por tejido
prostático de otros centros, se detecta con AUROC ≈ 0,91 en los mejores
métodos (ViM 0,912; RMDS 0,912; MDS 0,911). El control de resolución,
formado por las mismas imágenes del test ID degradadas a 128 px, alcanza
entre 0,753 y 0,789 en esos métodos. Al ser claramente inferior, la
detección de PANDA no se explica únicamente por la pérdida de nitidez:
el cambio de centro contribuye de forma propia. STOOD-X es la excepción,
con valores prácticamente iguales en ambos conjuntos (0,895 y 0,894).
Desde el punto de vista clínico, estos resultados implican que un
sistema validado en un hospital rechazaría una fracción importante de
imágenes de otro centro aun siendo del mismo órgano, lo que exigiría
recalibrar el detector antes de su despliegue.

\textbf{Far-OOD.} En DTD, STOOD-X obtiene el mejor resultado (AUROC
0,991; FPR@95 0,028). Llama la atención que, para ViM, MDS y KNN, el
far-OOD no resulta más fácil que el near-OOD (0,936--0,959 frente a
0,956--0,969). Esto refuerza la idea de que estos métodos son
especialmente sensibles a las características de adquisición de las
imágenes histológicas más que a su contenido semántico.

\textbf{Métodos basados en logits y en modificación de activaciones.}
MSP supera a MaxLogit y a Energy en los nueve conjuntos, al contrario de
lo habitual en la literatura. Una explicación plausible es el uso de
label smoothing durante el entrenamiento, que comprime la escala de los
logits; comprobarlo requeriría entrenar el modelo sin esta técnica. Los
métodos que modifican activaciones, diseñados para redes con
activaciones ReLU y BatchNorm, no se trasladan bien a un Swin
Transformer con LayerNorm. ReAct mejora a Energy, pero no a MSP. DICE
queda por debajo de Energy en ocho de los nueve conjuntos, con un AUROC
inferior al azar en DTD (0,436). ASH colapsa por completo: al
redistribuir la suma de las activaciones, que es negativa en el 27,5 \%
de las imágenes, invierte los logits y reduce la exactitud ID al 68,3
\%. Por último, SHE con producto escalar resulta irregular (0,795),
mientras que sus variantes con distancia coseno (0,938) y euclídea
(0,942) mejoran notablemente, lo que apunta de nuevo a la
incompatibilidad entre las características con signo del Swin y los
supuestos del método.

\textbf{Conclusión.} Ningún método domina en todas las métricas y
conjuntos. ViM ofrece el mejor equilibrio global en AUROC; RMDS, el
mejor comportamiento con umbrales exigentes, aunque con un fallo severo
en LYM; y STOOD-X destaca en far-OOD y en varias clases de NCT-CRC, pero
es el más sensible a la resolución. El hallazgo más relevante es que los
mejores detectores responden sobre todo a las condiciones de
adquisición, lo que condiciona su uso en un entorno multicéntrico.

\subsubsection{2.3.5. STOOD-X}\label{stood-x}

\begin{quote}
Se trata de una metodología post hoc de detección OOD en dos etapas: una
primera de detección, basada en un test estadístico no paramétrico sobre
distancias en el espacio de características, y una segunda de
explicación, que muestra al usuario los vecinos más cercanos y las
características más relevantes compartidas con ellos.

Sea x una imagen cuya clase predicha por el modelo es c, y sea X\_set el
conjunto de características de entrenamiento de esa clase. Se denota
X\_k = \{x\_1, \ldots, x\_K\} el conjunto de los K vecinos más cercanos
de x en X\_set según una distancia d, y, para cada vecino x\_i, X\_kⁱ el
conjunto de sus K vecinos más cercanos en X\_set \textbackslash{}
\{x\_i\}. La hipótesis del método es que, si x pertenece a la
distribución de entrenamiento, las distancias D(x, X\_k) y D(x\_i,
X\_kⁱ) siguen la misma distribución; si x es OOD, las primeras serán
mayores. Para contrastarlo se aplica el test unilateral de
Wilcoxon-Mann-Whitney entre ambos conjuntos de distancias, con un total
de M = K + K² distancias, y el p-valor resultante se utiliza como
puntuación: valores altos indican pertenencia a la distribución.

\textbf{Implementación.} El método se ha implementado siguiendo la
descripción del artículo.Se realiza un único test por imagen entre las K
distancias de x y las K² distancias de sus vecinos mediante
scipy.stats.mannwhitneyu (Virtanen et al., 2020).

\textbf{Configuración.} Se emplean la distancia coseno, K = 500 (el
valor seleccionado en el artículo), el conjunto de entrenamiento ID de
la clase predicha como referencia y dos configuraciones de
características: el 100 \% y el 87,5 \% más presente. Dado que el
criterio de relevancia del artículo se basa en CRP, que no es
directamente aplicable a un Swin Transformer, la selección del 87,5 \%
se realiza mediante la presencia de cada característica (suma de sus
valores absolutos en la clase), criterio utilizado en la implementación
pública. Ningún hiperparámetro se ha ajustado con datos OOD.

\textbf{Puntuación.}Como puntuación de detección se utiliza el
estadístico tipificado del test de Mann-Whitney, z, obtenido mediante la
aproximación normal del estadístico U con corrección por empates y por
continuidad. Dado que p = 1 − Φ(z) es una función monótona decreciente
de z, la puntuación −z induce la misma ordenación que el p-valor
propuesto en STOOD-X, pero evita el subdesbordamiento numérico que se
produce cuando p \textless{} 10⁻³⁰⁸. Cuando todas las distancias de x
superan a las de sus vecinos (separación completa), el estadístico
alcanza su valor máximo y varias imágenes quedan empatadas. Esto afecta
al 66 \% de TUM + NORM pero solo al 0,6 \% del test ID; su efecto está
acotado en ±0,003 en AUROC, y no altera FPR@95 ni CCR@FPR5.

\textbf{Resultados.} En el near-OOD principal (TUM + NORM), STOOD-X
obtiene un AUROC de 0,953, estadísticamente indistinguible del de RMDS
(0,954; p = 0,75 en el bootstrap pareado), aunque con un FPR@95
claramente mayor (0,349 frente a 0,128). Es el mejor método en NORM
(AUROC 0,990), STR (0,978), MUS (0,968) y DTD (0,991). En PANDA alcanza
0,895, prácticamente igual que en el control de resolución (0,894), lo
que indica que su detección del cambio de centro se explica casi por
completo por la diferencia de resolución. Las configuraciones con el 100
\% y el 87,5 \% de las características son indistinguibles.

\textbf{Análisis.}
\end{quote}

\begin{itemize}
\item
  \emph{Calibración.} Aunque el método propone el p-valor como medida de
  pertenencia, en el conjunto de validación ID el 82 \% de las imágenes
  obtiene p \textless{} 0,05, con una distribución bimodal. Por tanto,
  el p-valor permite ordenar las imágenes, pero no puede interpretarse
  como nivel de significación. La causa probable es que las K²
  distancias de referencia no son independientes, ya que los vecinos
  comparten vecinos. La calibración conforme sobre el conjunto de
  validación (sección X) recupera una tasa de falsas alarmas controlada.
\item
  \emph{Sensibilidad a K.} En un análisis de sensibilidad sobre un
  subconjunto (tabla X), el rendimiento mejora hasta K ≈ 72--288 y se
  degrada con K = 500. La degradación se concentra en las imágenes
  predichas en clases con pocas imágenes de entrenamiento
  (Gleason\_5\_comedonecrosis, 399; Inflamatorio, 568; No\_Cancer, 868),
  donde K se aproxima al tamaño de la clase y se pierde la localidad del
  vecindario. En particular, el 10 \% de las imágenes de TUM + NORM se
  predice como Gleason\_5\_comedonecrosis. Al excluir esas clases, la
  caída desaparece. Este análisis no se ha utilizado para elegir K, que
  se mantiene en el valor del artículo; una variante que adapta K al
  tamaño de la clase (K\_c = mín(500, ⌊n\_c/2⌋)) no mejoró los
  resultados.
\end{itemize}

\begin{quote}
\textbf{Explicaciones por ejemplos.} Siguiendo la segunda etapa del
método, la siguiente figura muestra imágenes representativas junto a sus
cinco vecinos más cercanos de entrenamiento. Las imágenes de colon
aceptadas como ID se predicen en su totalidad como
Gleason\_5\_comedonecrosis, lo que ilustra visualmente el efecto
descrito.

En cada figura, la primera columna es la imagen evaluada y las cinco
siguientes, sus vecinos de entrenamiento, con la clase predicha, la
puntuación (p y z) y las distancias en los títulos.

\includegraphics[width=6.26772in,height=5.80556in]{notas/media/image19.png}
\end{quote}

\textbf{Frases para usar en el TFG}

\emph{``Ambos conjuntos producen un cambio de distribución respecto al
conjunto ID, pero de naturaleza distinta: PANDA corresponde a un
desplazamiento de covariables, mientras que DTD introduce un cambio
semántico.''}

\emph{la detección OOD protege frente a imágenes raras o desconocidas,
pero no frente a cambios en el criterio diagnóstico. Para eso hace falta
revisar y reetiquetar los datos periódicamente.}

\includegraphics[width=6.26772in,height=2.02778in]{notas/media/image8.png}

\emph{Citar esta imagen y justificar:}

\emph{. En la memoria puedes usar esta figura (citándola) para situar tu
trabajo en (d) y justificar que incluyes PANDA porque, en un contexto
clínico, también interesa avisar de cambios de apariencia que degradan
el clasificador. Tu propio control lo demuestra: la exactitud cae del
96,87 \% al 90,77 \%.}

\emph{Cierta literatura previa clasifica las anomalías en tres tipos:
anomalías puntuales, anomalías condicionales o contextuales y anomalías
grupales o colectivas (Pang et al., 2020; Chalapathy y Chawla, 2019;
Ruff et al., 2021). En este estudio, nos centramos principalmente en la
detección de anomalías puntuales debido a su popularidad en aplicaciones
prácticas y a su idoneidad para esclarecer las similitudes y diferencias
entre las subtareas.}

\emph{Esto te sirve de respaldo teórico: puedes citar este párrafo, o
directamente a Fawcett (2006) para la curva ROC, al explicar tus
métricas.}

\emph{Formalmente, OSR aborda el caso en el que las muestras OOD
experimentan un cambio semántico durante las pruebas, es decir, P(Y) ≠
P′(Y). El objetivo de OSR es similar al de la detección de conjuntos
abiertos multiclase (ND); la única diferencia radica en que OSR
requiere, además, una clasificación precisa de las muestras ID de P(Y).
(CITAR)}

\textbf{Evaluación OSR TFG:}

Evaluación: Al igual que en AD y ND, las métricas para OSR incluyen
puntuaciones F (*F-scores*), AUROC y AUPR. Asimismo, el rendimiento de
la clasificación se evalúa mediante la precisión estándar para ID. Si
bien las métricas mencionadas evalúan de forma independiente las
capacidades de detección de novedades y de clasificación ID, algunos
trabajos proponen criterios para una evaluación conjunta, como CCR@FPRx
(Dhamija et al., 2018), que calcula la tasa de recuperación (*recall*)
por clase cuando se alcanza una determinada tasa de falsos positivos
(FPR) igual a *x* (por ejemplo, 10⁻¹).

\textbf{Evaluación OOD}

Evaluación Además de las puntuaciones F, AUROC y AUPR, otra métrica

comúnmente utilizada es FPR@TPRx, que mide la FPR cuando la TPR es x

(por ejemplo, 0,95). Algunos trabajos también utilizan una métrica
alternativa,

TNR@TPRx, que es equivalente a 1 - FPR@TPRx. La detección de OOD

también afecta al rendimiento de la clasificación de ID.

Aunque en clasificación general la detección OOD suele restringirse al
cambio semántico, Yang et al. (2024) señalan que en tareas de alto
riesgo, como la clasificación de imágenes médicas (Zimmerer et al.,
2022), se espera que los modelos sean conservadores y rechacen cualquier
muestra que se desvíe de la distribución de entrenamiento. Siguiendo
este criterio, en este trabajo se considera OOD tanto el cambio
semántico (DTD) como el desplazamiento de covariables debido al cambio
de centro (PANDA).''

\textbf{Añadir esta parte de Inferencia Estadística:}

Predicción Conforme La predicción conforme (CP) se erige como un marco
estadístico

robusto en el aprendizaje automático, diseñado principalmente para
proporcionar

medidas de confianza para las predicciones (Shafer y Vovk, 2008;
Angelopoulos y

Bates, 2021).

\emph{No se han incluido métodos basados en gradientes, como ODIN (Liang
et al., 2018) o GradNorm (Huang et al., 2021), por dos motivos. En
primer lugar, requieren una pasada hacia atrás por cada imagen, lo que
incrementa notablemente el coste de inferencia y dificulta su despliegue
en un entorno clínico. En segundo lugar, Igoe et al. (2022) indican que
la eficacia de estos métodos no se debe tanto a la información del
gradiente como a la magnitud de las características y a la distribución
de salida, aspectos que ya quedan cubiertos por los métodos basados en
características y en logits evaluados en este trabajo.''}

\emph{``Dentro de los métodos basados en densidad se distinguen dos
enfoques. El modelado mediante gaussianas condicionales por clase en el
espacio de características (Lee et al., 2018b) se ha incluido en la
evaluación a través de MDS y RMDS. En cambio, no se han considerado los
modelos generativos sobre el espacio de entrada, como los basados en
flujos, por tres motivos: su entrenamiento resulta costoso y difícil de
optimizar sobre imágenes de alta resolución; su verosimilitud presenta
un fuerte sesgo hacia la complejidad de la entrada (Serrà et al., 2020),
lo que en histopatología podría favorecer parches de baja complejidad,
como regiones de fondo o estroma; y su rendimiento suele quedar por
detrás de los enfoques basados en clasificación (Yang et al., 2024).''}

\emph{Yang et al. (2024) identifican la explicabilidad como uno de los
retos abiertos en la detección OOD. Este trabajo aborda dicho aspecto
desde dos perspectivas complementarias: mediante STOOD-X, que
proporciona explicaciones de las decisiones de detección, y mediante
técnicas de relevancia por conceptos (CRP), que permiten identificar y
eliminar características sesgadas del clasificador y analizar su efecto
sobre la fiabilidad del sistema.''}

\emph{Como puntuación de detección se utiliza el estadístico tipificado
del test de Mann-Whitney, z, obtenido mediante la aproximación normal
del estadístico U con corrección por empates y por continuidad. Dado que
p = 1 − Φ(z) es una función monótona decreciente de z, la puntuación −z
induce la misma ordenación que el p-valor propuesto en STOOD-X, pero
evita el subdesbordamiento numérico que se produce cuando p \textless{}
10⁻³⁰⁸.}

\emph{Hablar del efecto Clever Hans en XAI, y de como los modelos
aprender relaciones espurias para la clasificación de imággnes}

\emph{La diferencia que tiene CRP frente a los método de explicabilidad
local, es que estos otros reflejan las zonas en las que se ha fijado el
modelo pero no son capaces de determinar la importancia de estas
características para la decisión final}

\emph{Utilizar el concepto de glocal CRP al igual que se usa en el paper
que describe CRP}

{[}1{]} Kingma, D.; Ba, J. Adam: A Method for Stochastic Optimization.
arXiv 2014, arXiv:1412.6980.

Paper para métodos

{[}2{]} Kandel, I., \& Castelli, M. (2020). How Deeply to Fine-Tune a
Convolutional Neural Network: A Case Study Using a Histopathology
Dataset. \emph{Applied Sciences}, \emph{10}(10), 3359.
https://doi.org/10.3390/app10103359

Resolver problema CUDA out of memory

{[}3{]}
\href{https://saturncloud.io/blog/what-is-the-cuda-out-of-memory-error-and-how-to-fix-it/}{\ul{https://saturncloud.io/blog/what-is-the-cuda-out-of-memory-error-and-how-to-fix-it/}}

Paper con la explicación OOD

{[}4{]} Yang, J., Zhou, K., Li, Y. \emph{et al.} Generalized
Out-of-Distribution Detection: A Survey. \emph{Int J Comput Vis}
\textbf{132}, 5635--5662 (2024).
\href{https://doi.org/10.1007/s11263-024-02117-4}{\ul{https://doi.org/10.1007/s11263-024-02117-4}}

Explicación Test de Wilcoxon Mann Whitney

https://cienciadedatos.net/documentos/17\_mann\%E2\%80\%93whitney\_u\_test
